\documentclass{article}
\usepackage[spanish]{babel}
\usepackage{multicol}
\usepackage{amsmath}
\usepackage{amsfonts}
\usepackage{enumerate}
\usepackage[utf8]{inputenc}
\usepackage{amssymb}
\usepackage{graphicx}
\usepackage{wasysym}
\usepackage{moodle}
\usepackage[usenames]{color}
\newcommand{\radio}{\ooalign{\hidewidth$\bullet$\hidewidth\cr$\ocircle$}}
\newcommand{\R}{$\mathbb{R}^3$ }
\newcommand{\wi}{\displaystyle \int}
\newtheorem{theorem}{Theorem}
\usepackage[left=2.00cm, right=2.00cm, top=2.00cm, bottom=2.00cm]{geometry}
\title{Parcial 2}
\author{Departamento de matemáticas}
\date{\today}

\begin{document}
	\maketitle
	\begin{quiz}{}
    	\begin{cloze}{}
    	Considere el P.V.I
    	\[\begin{cases}\arctan(x)y'(t)-5\sin(y(t)) = 1, \ 1\leq t \leq 2 \\
    	y(1) = 4\end{cases}\]
        \begin{numerical}[points=15]{}
        	Tomando del m\'etodo de Euler modificado con $h=\frac{1}{10}$ se obtiene (con tres cifras decimales) un valor $y(1.6)\approx$\item[tolerance=0.003] 3.378
        \end{numerical}
        \begin{numerical}[points=18]{}
        	Ahora si se emplea el m\'etodo RK4 con $h=\frac{1}{10}$ se obtiene (con tres cifras decimales) un valor $y(1.6)\approx$\item[tolerance=0.003] 3.371
        \end{numerical}
        \end{cloze}
        \begin{cloze}{}
        	Considere el P.V.I
        	\[
        	\begin{cases}
        		e^{-t} y'(t) - 3\cos(y(t)) = t, \quad 1 \leq t \leq 2 \\
        		y(1) = 2
        	\end{cases}
        	\]
        	\begin{numerical}[points=15]{}
        		Tomando el m\'etodo de Euler modificado con $h=\frac{1}{10}$ se obtiene (con tres cifras decimales) un valor $y(1.6)\approx$ \item[tolerance=0.003] 2.104
        	\end{numerical}
        	\begin{numerical}[points=18]{}
        		Ahora si se emplea el m\'etodo RK4 con $h=\frac{1}{10}$ se obtiene (con tres cifras decimales) un valor $y(1.6)\approx$ \item[tolerance=0.003] 2.101
        	\end{numerical}
        \end{cloze}
    \begin{cloze}{}
    	Considere el P.V.I
    	\[
    	\begin{cases}
    		\sqrt{t+1} \, y'(t) - 2\sin(y(t)) = t^2, \quad 0 \leq t \leq 1 \\
    		y(0) = 3
    	\end{cases}
    	\]
    	\begin{numerical}[points=15]{}
    		Con el m\'etodo de Euler modificado con $h=\frac{1}{10}$ se obtiene $y(0.6)\approx$ \item[tolerance=0.003] 3.14
    	\end{numerical}
    	\begin{numerical}[points=18]{}
    		Con el m\'etodo RK4 con $h=\frac{1}{10}$ se obtiene $y(0.6)\approx$ \item[tolerance=0.003] 3.14
    	\end{numerical}
    \end{cloze}
    	\begin{cloze}{}
    		Considere el P.V.F:
    		\[\begin{cases}e^{-x}y''(x) + xy'(x) - \arctan(x)y(x) = 1+4x \\ 
    			y(1)=0\\ 
    			y(4) = 4
    		\end{cases}\]
    		\begin{numerical}[points=15]{}
    			Usando el m\'etodo de diferencias finitas con $h=\frac{3}{100}$ obtenemos un valor (con tres cifras decimales) para $y(1.3)\approx$\item[tolerance=0.003] -3.979
    		\end{numerical}
    	    \begin{numerical}[points=18]{}
    	    	Si usamos el m\'etodo del disparo lineal con $h=\frac{3}{100}$ entonces un valor aproximado(con tres cifras) para $y$ en $1.6$ es\item[tolerance=0.003] -15.533
    	    \end{numerical}
    	\end{cloze}
    \begin{cloze}{}
    	Considere el P.V.F:
    	\[
    	\begin{cases}
    		e^{-2x}y''(x) + x^2 y'(x) - \arctan(x)y(x) = 2 + 3x \\
    		y(0) = 1 \\
    		y(3) = 5
    	\end{cases}
    	\]
    	\begin{numerical}[points=15]{}
    		Usando el m\'etodo de diferencias finitas con $h=\frac{3}{100}$ obtenemos un valor (con tres cifras decimales) para $y(0.9)\approx$ \item[tolerance=0.003] -1.031
    	\end{numerical}
    	\begin{numerical}[points=18]{}
    		Si usamos el m\'etodo del disparo lineal con $h=\frac{3}{100}$ entonces un valor aproximado (con tres cifras) para $y$ en $1.5$ es \item[tolerance=0.003] -2.062
    	\end{numerical}
    \end{cloze}
        \begin{cloze}{}
        	Considere el P.V.I
        	\[\begin{cases}
        		y'''(x) = 5y'(x) + \arctan(x)e^xy(x) + 4, \ 1\leq x \leq 3 \\
        		y(1) = 1 \\
        		y'(1) = 1 \\
        		y''(1) = 1 \\
        	\end{cases}\]
    \begin{numerical}[points=11]{}
    	Si se aproxima la soluci\'on del P.V.I con $h=\frac{1}{10}$ obtenemos aproximaciones (con tres cifras) para $y(1.5)\approx$\item[tolerance=0.003] 1.906
    \end{numerical}
    \begin{numerical}[points=11]{}
    	para $y'(1.5)\approx$\item[tolerance=0.003] 3.332
    \end{numerical}
    \begin{numerical}[points=11]{}
    	y para $y''(1.5) \approx$\item[tolerance=0.003] 9.768
    \end{numerical}
        \end{cloze}
    \begin{cloze}{}
    	Considere el P.V.I
    	\[
    	\begin{cases}
    		y'''(x) = 4y'(x) + \sin(x)\, e^x\, y(x) + 3, \quad 0 \leq x \leq 2 \\\\
    		y(0) = 2 \\\\
    		y'(0) = -1 \\\\
    		y''(0) = 0 \\
    	\end{cases}
    	\]
    	\begin{numerical}[points=11]{}
    		Si se aproxima la soluci\'on del P.V.I con $h=\frac{1}{10}$ obtenemos aproximaciones (con tres cifras) para $y(1)\approx$ \item[tolerance=0.003] 0.895
    	\end{numerical}
    	\begin{numerical}[points=11]{}
    		para $y'(1)\approx$ \item[tolerance=0.003] -1.275
    	\end{numerical}
    	\begin{numerical}[points=11]{}
    		y para $y''(1)\approx$ \item[tolerance=0.003] -0.448
    	\end{numerical}
    \end{cloze}
    \end{quiz}
\end{document}